% Phase-2 fix designs, group G3 (first-order chain): R-007, R-008, R-009, R-010, R-014.
% Line numbers and typewriter labels are anchors into source/formal-logic-reference.tex
% (v1.0, 14780 lines).  Plain LaTeX; compilable under amsmath/amssymb.
% Set in \small so that each design occupies at most one page.
\begingroup\small

\subsection*{R-007 --- Prenex normal form: missing lemma chain (v1.0 lines 4396--4446)}

\paragraph{Defect.}
The prenex theorem's proof (4424--4446) opens with unproved global operations:
``move negations inward using De Morgan equivalences and quantifier duality,
and alpha-rename bound variables so that the formula is rectified''
(4425--4427).  No first-order negation-normal-form lemma exists (the
propositional theorem, 1899--1917, covers no quantifiers and its termination
measure counts only connectives below a negation); no rectified-variant
existence result exists (no renaming recursion is ever defined); and
``quantifier duality'' resolves to \texttt{lem:quantifier-duality} (3814), a
derivability lemma of \(\mathsf{ND}_{\mathrm{FOL}}\), where a semantic
equivalence is required.  The extraction induction then rests on
``Rectification guarantees the side condition'' (4441) without any stated
invariant.

\paragraph{Repair.}
Insert Lemmas A--C between the definitions (4399--4418) and the theorem
(4420), rewrite the theorem's proof as their composition, and attach the
missing label \texttt{cor:alpha-invariance} at 3203; \(c\) is the complexity
of 2875 (counting quantifiers, so the measure covers quantifier prefixes) and
\(\equiv\) is logical equivalence (3256).

\emph{Lemma A (\texttt{lem:fol-nnf}, negation normal form).}  Define
\(\operatorname{nnf},\operatorname{nnf}':\mathrm{Fm}(\mathcal{L})\to\mathrm{Fm}(\mathcal{L})\)
by simultaneous recursion: \(\operatorname{nnf}\) fixes atomic formulas and
\(\bot\), commutes with \(\wedge,\vee,\forall,\exists\), and satisfies
\(\operatorname{nnf}(\neg\varphi)=\operatorname{nnf}'(\varphi)\),
\(\operatorname{nnf}(\varphi\to\psi)=\operatorname{nnf}'(\varphi)\vee\operatorname{nnf}(\psi)\);
and, for \(\alpha\) atomic,
\(\operatorname{nnf}'(\alpha)=\neg\alpha\),
\(\operatorname{nnf}'(\bot)=\neg\bot\),
\(\operatorname{nnf}'(\neg\varphi)=\operatorname{nnf}(\varphi)\),
\(\operatorname{nnf}'(\varphi\wedge\psi)=\operatorname{nnf}'(\varphi)\vee\operatorname{nnf}'(\psi)\),
\(\operatorname{nnf}'(\varphi\vee\psi)=\operatorname{nnf}'(\varphi)\wedge\operatorname{nnf}'(\psi)\),
\(\operatorname{nnf}'(\varphi\to\psi)=\operatorname{nnf}(\varphi)\wedge\operatorname{nnf}'(\psi)\),
\(\operatorname{nnf}'(\forall x\,\varphi)=\exists x\,\operatorname{nnf}'(\varphi)\),
\(\operatorname{nnf}'(\exists x\,\varphi)=\forall x\,\operatorname{nnf}'(\varphi)\)
(every call is on an immediate subformula, so \(c\) strictly decreases and the
recursion is legitimate).  Claims, by simultaneous strong induction on
\(c(\varphi)\): (i) both values are in negation normal form (definition
4399--4408; \(\neg\bot\) is admitted there); (ii)
\(\operatorname{nnf}(\varphi)\equiv\varphi\) and
\(\operatorname{nnf}'(\varphi)\equiv\neg\varphi\), each case checked against
the satisfaction clauses (3111--3145): the \(\operatorname{nnf}'\) quantifier
cases are the pointwise computations
\(\mathcal{M},s\models\neg\forall x\,\varphi\) iff some \(a\) has
\(\mathcal{M},s[x\mapsto a]\models\neg\varphi\) iff
\(\mathcal{M},s\models\exists x\,\neg\varphi\), dually for \(\exists\) ---
not citations of \texttt{lem:quantifier-duality}; (iii) \(\mathrm{FV}\) and
\(\mathrm{Var}\) of both values equal those of \(\varphi\).

\emph{Lemma B (\texttt{lem:rectified-variant}, rectified variant).}  For
\(\varphi\) and finite \(V\supseteq\mathrm{Var}(\varphi)\) define
\(r(\varphi,V)\): atomic and \(\bot\) unchanged; \(r(\neg\varphi,V)=\neg
r(\varphi,V)\); \(r(\varphi\circ\psi,V)=r(\varphi,V)\circ r(\psi,\,V\cup
\mathrm{Var}(r(\varphi,V)))\) for \(\circ\in\{\wedge,\vee,\to\}\);
\(r(Qx\,\varphi,V)=Qz\,r(\varphi[z/x],\,V\cup\{z\})\) with \(z\) least outside
\(V\).  Well-foundedness is the 2951--2958 pattern:
\(z\notin\mathrm{Var}(\varphi)\cup\{x\}\), so no binder of \(\varphi\) equals
\(z\), the substitution runs on the first two quantifier clauses only (textual
replacement, so
\(\mathrm{Var}(\varphi[z/x])\subseteq\mathrm{Var}(\varphi)\cup\{z\}\)), and
\(c(\varphi[z/x])=c(\varphi)<c(Qx\,\varphi)\) by
\texttt{lem:substitution-complexity} (2965).  Claims, by strong induction on
\(c(\varphi)\) under the invariant \(\mathrm{Var}(\varphi)\subseteq V\):
(i) \(r(\varphi,V)\equiv_{\alpha}\varphi\) --- connective cases by congruence,
quantifier case by the generating clause \(Qx\,\varphi\equiv_\alpha
Qz\,(\varphi[z/x])\) (2984, licensed by \(z\notin\mathrm{Var}(\varphi)\)), the
\(Q\)-congruence, and transitivity --- hence \(r(\varphi,V)\equiv\varphi\) by
\texttt{cor:alpha-invariance}; (ii) \(r(\varphi,V)\) is rectified: its binders
are pairwise distinct and outside \(V\) (the binary clause forbids the left
output's variables to the right recursion), while
\(\mathrm{FV}(r(\varphi,V))=\mathrm{FV}(\varphi)\subseteq V\) by (iii);
(iii) alpha-equivalent formulas have equal free-variable sets: the
equal-\(\mathrm{FV}\) relation is a congruence equivalence containing each
generating pair, since both cases of clause~(b) of
\texttt{lem:substitution-fv} (added by R-010) give
\(\mathrm{FV}(Qy\,(\varphi[y/x]))=\mathrm{FV}(Qx\,\varphi)\), so leastness of
\(\equiv_\alpha\) applies; (iv) \(r\) preserves the constructor tree, hence
negation normal form.

\emph{Lemma C (\texttt{lem:prenex-extraction}, extraction with invariant).}
For every rectified \(\varphi\) in negation normal form there is
\(p(\varphi)\) with: (i) \(p(\varphi)\equiv\varphi\); (ii) \(p(\varphi)\)
rectified and prenex with quantifier-free negation-normal matrix; (iii)
\emph{invariant:} \(\mathrm{FV}\) and \(\mathrm{Var}\) unchanged, and the
prefix consists of exactly the quantifier occurrences of \(\varphi\) (same
quantifier, same variable) --- no new free or bound variables.  Route:
structural induction on the negation-normal generation (subformulas of
rectified formulas are rectified); \(p\) is the identity on base formulas and
commutes with \(Qx\).  Binary case \(\alpha\circ\beta\),
\(\circ\in\{\wedge,\vee\}\): with
\(p(\alpha)=Q_1y_1\cdots Q_ky_k\,\alpha_0\),
\(p(\beta)=R_1z_1\cdots R_mz_m\,\beta_0\), set
\(p(\alpha\circ\beta)=Q_1y_1\cdots Q_ky_k\,R_1z_1\cdots R_mz_m\,
(\alpha_0\circ\beta_0)\), proved \(\equiv\alpha\circ\beta\) by an inner
induction on \(k+m\) from two recorded facts: (C1)
\((Qy\,\pi)\circ\chi\equiv Qy\,(\pi\circ\chi)\) and
\(\chi\circ(Qy\,\pi)\equiv Qy\,(\chi\circ\pi)\) whenever
\(y\notin\mathrm{FV}(\chi)\), proved from the satisfaction clauses and
\texttt{lem:fol-coincidence} (3148), displaying the \(\forall\)/\(\vee\) case
(the classical distribution of \(\forall\) over a disjunct constant in the
quantified variable); (C2) \(\equiv\) is a congruence for \(\circ\) and
\(Qx\) (immediate from the satisfaction clauses).  The side conditions hold
because \(y_i\notin\mathrm{FV}(p(\beta))=\mathrm{FV}(\beta)\) (\(y_i\) is
bound in the rectified \(\alpha\circ\beta\)) and
\(z_j\notin\mathrm{FV}(\alpha_0)\subseteq
\mathrm{FV}(\alpha)\cup\{y_1,\ldots,y_k\}\) (rectifiedness, binder
distinctness); the output is rectified and satisfies the invariant since
nothing is renamed or introduced.

\emph{Rewritten theorem proof} (replaces 4424--4446; the statement gains label
\texttt{thm:prenex-normal-form}): by Lemmas A--C,
\(p\bigl(r(\operatorname{nnf}(\varphi),\mathrm{Var}(\operatorname{nnf}(\varphi)))\bigr)
\equiv\varphi\) is rectified, prenex, and in the same signature; all four
results then carry \textit{Status: complete proof}.

\paragraph{Acceptance.}
(1) The rewritten proof cites only Lemmas A--C and
\texttt{cor:alpha-invariance}; ``quantifier duality'' no longer occurs in
4396--4446.  (2) Each lemma names its induction measure; Lemma B's renaming
recursion cites \texttt{lem:substitution-complexity} for its non-structural
call on \(\varphi[z/x]\).  (3) Lemma C states the no-new-variables invariant and every extraction
step's side condition is discharged from it plus rectifiedness.

\subsection*{R-008 --- Skolemization: explicit prefix induction replacing the ``quantifier game'' (v1.0 lines 4480--4512)}

\paragraph{Defect.}
Both directions of \texttt{thm:skolem-equisatisfiable} (4492--4507) are prose:
``Induction along the prefix then shows\ldots'' states no induction hypothesis,
and the converse appeals to ``the original quantifier game'' (4501) --- a notion
defined nowhere in the document.  The Axiom of Choice is hypothesized for the
whole theorem (4482) although only the expansion direction uses it.

\paragraph{Repair.}
Fix \(\sigma=Q_1x_1\cdots Q_nx_n\,\theta\) rectified prenex.  For \(0\le i\le
n\) let \(\sigma_i=Q_{i+1}x_{i+1}\cdots Q_nx_n\,\theta\) (so
\(\sigma_0=\sigma\), \(\sigma_n=\theta\), and
\(\mathrm{FV}(\sigma_i)\subseteq\{x_1,\ldots,x_i\}\)); let \(q(i)\) count the
universal positions among the first \(i\), let \(\bar u(j)\) list the universal
variables before position \(j\), and for existential \(j\) let
\(t_j=f_j(\bar u(j))\) be its Skolem term, so
\(\operatorname{Sk}(\sigma)=\forall u_1\cdots\forall u_p\,\theta^{*}\) with
\(\theta^{*}=\theta[t_j/x_j]\) simultaneously for all existential \(j\)
(\(\theta\) is quantifier-free, so this is plain replacement; since no \(t_j\)
contains an existential variable, the simultaneous substitution equals the
composition of the single substitutions in any order).  Given a structure and a
tuple \(\bar a\) of values for the universal variables among
\(x_1,\ldots,x_i\), define the \emph{Skolem-term assignment} \(s_{\bar a}\):
universal \(x_j\) gets its entry of \(\bar a\); existential \(x_j\) gets the
value \(v_j\) of \(t_j\) computed from the entries for \(\bar u(j)\); values of
other variables are irrelevant by \texttt{lem:fol-coincidence}.

\emph{Restated theorem} (replaces 4480--4491): let \(\sigma\) be a rectified
prenex \(\mathcal{L}\)-sentence.  (1) \emph{(Axiom of Choice.)}  Every
\(\mathcal{L}\)-structure satisfying \(\sigma\) has an expansion to the Skolem
signature satisfying \(\operatorname{Sk}(\sigma)\).  (2) \emph{(No choice.)}
The \(\mathcal{L}\)-reduct of every model of \(\operatorname{Sk}(\sigma)\)
satisfies \(\sigma\).  Hence, assuming the Axiom of Choice, \(\sigma\) has a
model iff \(\operatorname{Sk}(\sigma)\) has one; the right-to-left implication
holds without choice.

\emph{Proof of (1).}  Let \(\mathcal{M}\models\sigma\).  By finite recursion on
\(i\) choose set-functions \(F_j\) (for existential \(j\le i\)) while proving
\[
 \mathrm{Inv}(i):\quad
 \mathcal{M},s_{\bar a}\models\sigma_i
 \quad\text{for every }\bar a\in|\mathcal{M}|^{\,q(i)},
\]
where \(s_{\bar a}\) computes \(v_j=F_j(\text{entries of }\bar a\text{ for }
\bar u(j))\).  \(\mathrm{Inv}(0)\) is the hypothesis.  Universal step: from
\(\mathcal{M},s_{\bar a}\models\forall x_{i+1}\,\sigma_{i+1}\), the
satisfaction clause gives \(\mathcal{M},s_{\bar a}[x_{i+1}\mapsto b]\models
\sigma_{i+1}\) for every \(b\), which is \(\mathrm{Inv}(i+1)\) for
\((\bar a,b)\).  Existential step: by \(\mathrm{Inv}(i)\) each witness set
\(W(\bar a)=\{b:\mathcal{M},s_{\bar a}[x_{i+1}\mapsto b]\models\sigma_{i+1}\}\)
is nonempty; apply the Axiom of Choice (definition at 651) to
\(\{\{\bar a\}\times W(\bar a):\bar a\in|\mathcal{M}|^{\,q(i)}\}\) to obtain
\(F_{i+1}\) with \(F_{i+1}(\bar a)\in W(\bar a)\) --- one application of choice
per existential quantifier, and nowhere else.  Define the expansion
\(\mathcal{M}^{*}\) by \(f_j^{\mathcal{M}^{*}}=F_j\).  From
\(\mathrm{Inv}(n)\): for every \(\bar a\in|\mathcal{M}|^{p}\), the term-value
clauses give \(t_j\) the value \(v_j\) under \([\bar u\mapsto\bar a]\) in
\(\mathcal{M}^{*}\), so iterating \texttt{thm:fol-substitution} once per
existential position turns \(\mathcal{M},s_{\bar a}\models\theta\) (equal to
\(\mathcal{M}^{*},s_{\bar a}\models\theta\) by \texttt{lem:reduct-invariance})
into \(\mathcal{M}^{*},[\bar u\mapsto\bar a]\models\theta^{*}\); the
\(p\)-fold universal clause yields
\(\mathcal{M}^{*}\models\operatorname{Sk}(\sigma)\).

\emph{Proof of (2).}  Let \(\mathcal{N}\models\operatorname{Sk}(\sigma)\) in
the expanded signature; now \(v_j\) is the value of the term \(t_j\) in
\(\mathcal{N}\).  By downward induction from \(i=n\) to \(i=0\) prove
\(\mathrm{Inv}'(i)\): \(\mathcal{N},s_{\bar a}\models\sigma_i\) for every
\(\bar a\in|\mathcal{N}|^{\,q(i)}\).  Base \(i=n\): instantiate the universal
prefix of \(\operatorname{Sk}(\sigma)\) at \(\bar a\) and reverse the same
iterated substitution-lemma computation as in (1).  Universal step
(\(Q_{i+1}=\forall\)): for every \(b\), \(\mathrm{Inv}'(i+1)\) at
\((\bar a,b)\) gives \(\mathcal{N},s_{\bar a}[x_{i+1}\mapsto b]\models
\sigma_{i+1}\) (the \(v_j\), \(j\le i\), depend only on \(\bar a\)); apply the
\(\forall\)-clause.  Existential step: \(\mathrm{Inv}'(i+1)\) at \(\bar a\)
exhibits the witness \(v_{i+1}\), and the \(\exists\)-clause gives
\(\mathrm{Inv}'(i)\).  Then \(\mathrm{Inv}'(0)\) and
\texttt{lem:reduct-invariance} yield
\(\mathcal{N}\mathbin{\upharpoonright}\mathcal{L}\models\sigma\).  No choice
principle is used.  The remark at 4509--4512 stays.  The theorem carries
\textit{Status: complete proof} afterwards.

\paragraph{Acceptance.}
(1) The string ``quantifier game'' occurs nowhere in the manuscript.  (2) Each
direction displays an induction statement indexed by prefix position with the
Skolem-term assignment defined, and every step is discharged by a named
satisfaction clause, \texttt{thm:fol-substitution},
\texttt{lem:fol-coincidence}, or \texttt{lem:reduct-invariance}.  (3) The Axiom
of Choice appears only in clause (1), invoked exactly once per existential
quantifier on a displayed family of witness sets; clause (2)'s proof invokes no
choice principle.

\subsection*{R-009 --- Transfinite Lindenbaum lemma for the Henkin construction (v1.0 lines 3934--3985)}

\paragraph{Defect.}
The second paragraph of the Henkin--Lindenbaum proof (3978--3985) says
``Well-order all formulas of \(\mathcal{L}^{H}\) and repeat the Lindenbaum
construction'' and claims maximal consistency ``by the same finite-derivation
argument as in \texttt{lem:prop-lindenbaum}.''  That lemma (2381--2405) is an
\(\omega\)-indexed recursion over a countable enumeration; for
\(\kappa\) possibly uncountable the construction needs transfinite recursion
with limit stages, which the propositional argument does not contain, and the
family whose union is taken is never defined.  The witness stage two paragraphs
earlier (3961--3976) does all of this explicitly via
\texttt{thm:transfinite-recursion}; the two halves of one proof sit at
different rigor levels.

\paragraph{Repair.}
Insert between 3932 and 3934:

\emph{Lemma (\texttt{lem:transfinite-lindenbaum}).}  Let \(\mathcal{L}'\) be a
signature, \(\lambda\) an ordinal, and
\(\langle\varphi_\alpha:\alpha<\lambda\rangle\) an enumeration of
\(\mathrm{Fm}(\mathcal{L}')\) in which every formula occurs at least once.
Then every set \(\Sigma\subseteq\mathrm{Fm}(\mathcal{L}')\) consistent for
\(\mathsf{ND}_{\mathrm{FOL}}\) is contained in a maximal consistent
\(\Delta\subseteq\mathrm{Fm}(\mathcal{L}')\).  (Choice-free given the
enumeration; the Axiom of Choice enters only when producing one.)

\emph{Proof route.}  By \texttt{thm:transfinite-recursion} (632) with domain
\(\lambda+1\) define \(\langle\Gamma_\alpha:\alpha\le\lambda\rangle\):
\[
 \Gamma_0=\Sigma,\qquad
 \Gamma_{\alpha+1}=
 \begin{cases}
   \Gamma_\alpha\cup\{\varphi_\alpha\},&\text{if this set is consistent},\\
   \Gamma_\alpha,&\text{otherwise},
 \end{cases}
 \qquad
 \Gamma_\beta=\bigcup_{\alpha<\beta}\Gamma_\alpha\ \text{ for limit }\beta.
\]
\emph{Chain claim:} \(\alpha\le\beta\le\lambda\) implies
\(\Gamma_\alpha\subseteq\Gamma_\beta\) (transfinite induction on \(\beta\)).
\emph{Consistency:} transfinite induction on \(\alpha\); zero by hypothesis;
successor by the case split; at a limit \(\beta\), a derivation of \(\bot\)
from \(\Gamma_\beta\) is finite, so its open assumptions form a finite
\(F\subseteq\Gamma_\beta\); each \(\psi\in F\) lies in some
\(\Gamma_{\alpha_\psi}\) with \(\alpha_\psi<\beta\), and with
\(\alpha^{*}=\max_{\psi\in F}\alpha_\psi<\beta\) if \(F\ne\varnothing\) (a
finite nonempty set of ordinals has a maximum) and \(\alpha^{*}=0<\beta\) if
\(F=\varnothing\), the chain claim locates \(F\subseteq\Gamma_{\alpha^{*}}\),
contradicting the induction hypothesis at \(\alpha^{*}\).  \emph{Maximality of}
\(\Delta=\Gamma_\lambda\): if \(\varphi\notin\Delta\), pick \(\alpha\) with
\(\varphi_\alpha=\varphi\); rejection at stage \(\alpha+1\) means
\(\Gamma_\alpha\cup\{\varphi\}\) derives \(\bot\), and since a derivation from
a set is a derivation from any superset (monotonicity, immediate from the
definition of derivability from assumptions),
\(\Delta\cup\{\varphi\}\) derives \(\bot\).

\emph{Rewrite of 3978--3985:} ``The set \(\mathrm{Fm}(\mathcal{L}^{H})\) has
cardinality exactly \(\kappa\), as established above from
\texttt{thm:infinite-cardinal-finite-coding} and the formulas
\(\exists x_0\,(x_0=d_\xi)\) (the count at 3949--3952 is retained); by the
cardinal-assignment corollary (662, under the Axiom of Choice already assumed)
fix a bijective enumeration
\(\langle\varphi_\alpha:\alpha<\kappa\rangle\) of it.  Apply
\texttt{lem:transfinite-lindenbaum} to \(\Sigma^{H}\) to obtain a maximal
consistent \(\Delta\supseteq\Sigma^{H}\supseteq\Sigma\).''  The existing
witness-property sentences (3981--3984) follow unchanged.  Status labels
afterwards: the new lemma and \texttt{lem:henkin-lindenbaum} both
\textit{Status: complete proof.}

\paragraph{Acceptance.}
(1) A lemma between \texttt{lem:new-constant-elimination} and
\texttt{lem:henkin-lindenbaum} displays the recursion with zero, successor, and
limit clauses and cites \texttt{thm:transfinite-recursion}.  (2) Its limit case
contains the chain claim and locates an arbitrary finite subset of a limit
stage inside a single earlier stage via the maximum of finitely many ordinals.
(3) Lines 3978--3981 no longer contain ``repeat the Lindenbaum construction''
or the appeal to \texttt{lem:prop-lindenbaum}; the paragraph applies the new
lemma to \(\Sigma^{H}\) under an explicitly fixed enumeration of order type
\(\kappa\).

\subsection*{R-010 --- Vacuous substitution is alpha-invariant: strong induction on complexity (v1.0 lines 3006--3016)}

\paragraph{Defect.}
The unnamed lemma at 3006 (``If \(x\notin\mathrm{FV}(\varphi)\), then
\(\varphi[t/x]\equiv_\alpha\varphi\)'') is proved ``by structural induction,''
but in the renaming case the recursion visits \(\psi[z/y]\), which is not a
subformula of \(Qy\,\psi\), so plain structural induction does not license the
step; the case is discharged by the bare assertion that renaming ``changes the
formula only within its alpha-equivalence class.''  The book's own template for
this situation is the simultaneous strong induction on complexity at 2951--2958
and its use at 3193--3199 of \texttt{lem:substitution-complexity}.

\paragraph{Repair.}
First insert, after the example ending at 3004 and before the lemma at 3006:

\emph{Lemma (\texttt{lem:substitution-fv}, free variables under
substitution).}  For every term \(r\): \(x\notin\mathrm{Var}(r)\) implies
\(r[t/x]=r\), and
\(\mathrm{Var}(r[t/x])\subseteq(\mathrm{Var}(r)\setminus\{x\})\cup\mathrm{Var}(t)\)
(induction on \(r\)).  For every formula \(\varphi\), term \(t\), and variables
\(x,y\):
\begin{itemize}
\item[(a)] \(\mathrm{FV}(\varphi[t/x])\subseteq
      (\mathrm{FV}(\varphi)\setminus\{x\})\cup\operatorname{Var}(t)\)
      (the book defines no \(\mathrm{FV}\) on terms; a term's variables
      are its \(\operatorname{Var}\), as in the \(\mathrm{FV}\) clauses for
      atomic formulas at 2862--2863);
\item[(b)] if \(y\notin\mathrm{Var}(\varphi)\), then
      \(\mathrm{FV}(\varphi[y/x])=(\mathrm{FV}(\varphi)\setminus\{x\})\cup\{y\}\)
      when \(x\in\mathrm{FV}(\varphi)\), and
      \(\mathrm{FV}(\varphi[y/x])=\mathrm{FV}(\varphi)\) when
      \(x\notin\mathrm{FV}(\varphi)\).
\end{itemize}
Route: strong induction on \(c(\varphi)\).  For (b) the renaming clause of the
substitution definition never fires (a binder \(w\) of a subformula lies in
\(\mathrm{Var}(\varphi)\), so \(w\ne y\)), and the substitution is textual
replacement.  For (a)'s renaming case
\((Qy\,\psi)[t/x]=Qz\,((\psi[z/y])[t/x])\): apply the induction hypothesis to
\(\psi\) with \([z/y]\) and to \(\psi[z/y]\) with \([t/x]\) --- the latter at
equal complexity \(c(\psi[z/y])=c(\psi)<c(Qy\,\psi)\), licensed by
\texttt{lem:substitution-complexity} (2965) --- then compute, using
\(z\notin\mathrm{Var}(t)\cup\{x\}\), that
\(\mathrm{FV}(Qz\,((\psi[z/y])[t/x]))\subseteq
(\mathrm{FV}(Qy\,\psi)\setminus\{x\})\cup\operatorname{Var}(t)\).

Then replace the proof at 3010--3016 by: \emph{Proceed by strong induction on
formula complexity \(c\) (2875).}  Atomic case: \(x\) occurs in no term, so
every term is unchanged by the term clause of \texttt{lem:substitution-fv} and
\(\varphi[t/x]=\varphi\); \(\bot\) likewise.  Connective cases:
\(x\notin\mathrm{FV}\) of each immediate subformula; induction hypotheses and
the congruence property of \(\equiv_\alpha\).  Quantifier \(Qy\,\psi\) with
\(y=x\): substitution stops, reflexivity.  With \(y\ne x\),
\(y\notin\mathrm{Var}(t)\): then \(x\notin\mathrm{FV}(\psi)\); the induction
hypothesis gives \(\psi[t/x]\equiv_\alpha\psi\), and the \(Q\)-congruence
concludes.  With \(y\ne x\), \(y\in\mathrm{Var}(t)\): the definition returns
\(Qz\,((\psi[z/y])[t/x])\) with \(z\) least outside
\(\mathrm{Var}(\psi)\cup\mathrm{Var}(t)\cup\{x,y\}\).  By clause (b) of
\texttt{lem:substitution-fv}, \(x\notin\mathrm{FV}(\psi[z/y])\) (as
\(x\notin\mathrm{FV}(\psi)\) and \(x\ne z\)).  By
\texttt{lem:substitution-complexity}, \(c(\psi[z/y])=c(\psi)<c(Qy\,\psi)\), so
the \emph{strong} induction hypothesis applies to \(\psi[z/y]\) and gives
\((\psi[z/y])[t/x]\equiv_\alpha\psi[z/y]\); the \(Q\)-congruence gives
\(Qz\,((\psi[z/y])[t/x])\equiv_\alpha Qz\,(\psi[z/y])\); the generating clause
(licensed since \(z\notin\mathrm{Var}(\psi)\)) gives
\(Qz\,(\psi[z/y])\equiv_\alpha Qy\,\psi\); transitivity concludes.  Landing:
new lemma at approximately 3005; rewritten proof in place at 3010--3016.  Both
results carry \textit{Status: complete proof} afterwards.  (Clause (b) is also
consumed by R-007's Lemma B.)

\paragraph{Acceptance.}
(1) The proof at 3010--3016 opens ``by strong induction on formula
complexity'' and cites \texttt{lem:substitution-complexity} exactly where the
induction hypothesis is applied to \(\psi[z/y]\).  (2) A stated and proved
\texttt{lem:substitution-fv} precedes the lemma and is cited for
\(x\notin\mathrm{FV}(\psi[z/y])\).  (3) No case applies an induction
hypothesis to a formula that is neither of strictly smaller complexity nor of
equal complexity justified by \texttt{lem:substitution-complexity}.

\subsection*{R-014 --- \(\mathsf{LK}\) completeness for arbitrary sequents (v1.0 lines 4165--4185, 2488--2507; soundness 3674--3687, 2295--2318)}

\paragraph{Defect.}
Soundness (3674) covers every derivable sequent: the calculus is
multi-succedent (rules at 2178--2234 and 3583--3618, satisfaction convention at
3579--3581), and the theorem quantifies over sequents
\(\Gamma\Rightarrow\Delta\).  Completeness (4165) is stated only ``for every
finite multiset \(\Gamma\) and formula \(\varphi\)'' --- single-succedent.  No
completeness statement covers sequents with empty or multi-formula succedents,
so the two theorems do not compose to a derivability-iff-validity equivalence
for the objects the calculus manipulates.  The propositional pair exhibits the
identical asymmetry --- \texttt{thm:prop-lk-soundness} (2295--2318) covers
every derivable \(\mathsf{LK}_{\mathrm{CPL}}\) sequent while completeness
(2488--2491) is again single-succedent --- and register row R-014 spans both
pairs (2488 vs 2295; 4165 vs 3674), so this design must repair both.

\paragraph{Repair.}
Attach \texttt{thm:fol-lk-soundness} to the theorem at 3674.  Replace the
statement at 4165--4170 by:

\emph{Theorem (\texttt{thm:fol-lk-completeness}).}  Assume the Axiom of Choice
in the ambient set theory.  For all finite multisets \(\Gamma,\Delta\) of
\(\mathcal{L}\)-formulas, the sequent \(\Gamma\Rightarrow\Delta\) is derivable
in \(\mathsf{LK}_{\mathrm{FOL}}\) iff for every \(\mathcal{L}\)-structure
\(\mathcal{M}\) and assignment \(s\): if \(\mathcal{M},s\models\gamma\) for
every \(\gamma\in\Gamma\), then \(\mathcal{M},s\models\delta\) for some
\(\delta\in\Delta\).

(The roadmap's propositional phrasing ``every valuation'' is this schema
instantiated for \(\mathsf{LK}_{\mathrm{CPL}}\); at first order the semantic
data are structure plus assignment, matching 3579--3581.)  Left to right is
\texttt{thm:fol-lk-soundness}.  Right to left, by reduction to
\texttt{thm:fol-completeness} (4128) plus the existing translation:
(1)~let \(\Delta=\delta_1,\ldots,\delta_m\) (\(m\ge 0\)) and let \(S\) be the
set of members of \(\Gamma\) together with
\(\neg\delta_1,\ldots,\neg\delta_m\); validity of the sequent makes \(S\)
unsatisfiable, i.e.\ \(S\models\bot\), since a structure and assignment
satisfying all of \(S\) would satisfy some \(\delta_j\) and its negation.
(2)~\texttt{thm:fol-completeness} gives
\(S\vdash_{\mathsf{ND}_{\mathrm{FOL}}}\bot\); the derivation's open assumptions
form a finite \(S_0\subseteq S\), say \(\gamma_1,\ldots,\gamma_k\in\Gamma\)
and \(\neg\delta_{j_1},\ldots,\neg\delta_{j_r}\); the translation of the
existing sketch (4171--4185) turns it into an
\(\mathsf{LK}_{\mathrm{FOL}}\)-derivation of
\(\gamma_1,\ldots,\gamma_k,\neg\delta_{j_1},\ldots,\neg\delta_{j_r}
\Rightarrow\bot\).
(3)~Cut with the initial \(\bot\)-sequent \(\bot\Rightarrow\;\) (instance of
\(\Gamma,\bot\Rightarrow\Delta\) with empty contexts, 2180) empties the
succedent.
(4)~For each \(j\), the sequent \(\;\Rightarrow\delta_j,\neg\delta_j\) is
derivable (identity sequent \(\delta_j\Rightarrow\delta_j\), then \(\neg R\));
cutting it against the current sequent, with cut formula \(\neg\delta_j\),
moves \(\delta_j\) into the succedent:
\[
 \frac{\;\Rightarrow\delta_j,\neg\delta_j
 \qquad \Gamma',\neg\delta_j\Rightarrow\Pi}
 {\Gamma'\Rightarrow\delta_j,\Pi}
 \qquad\text{yielding, after }r\text{ cuts,}\qquad
 \gamma_1,\ldots,\gamma_k\Rightarrow\delta_{j_1},\ldots,\delta_{j_r}.
\]
(5)~Weakening (\(W_L,W_R\)) adds the unused members of \(\Gamma,\Delta\) and
contraction (\(C_L,C_R\)) merges duplicates, giving exactly
\(\Gamma\Rightarrow\Delta\) as multisets.  The case \(m=0\) runs with no
step-(4) cuts and asserts that no structure and assignment satisfies all of
\(\Gamma\); empty \(\Gamma\) is likewise covered.  The route uses cut, which is
primitive in the displayed calculus; no cut-free claim is made, and extending
the sketch's translation induction to multi-succedent targets is unnecessary
since the reduction reuses it verbatim.  Landing: statement replaced at
4165--4170; steps (1)--(5) appended to the proof at 4171--4185, whose
translation subroutine is retained.  Status afterwards: \textit{Status: proof
sketch} while the rule-by-rule translation trees remain omitted; the theorem
inherits \textit{Status: complete proof} only if a separate work order supplies
them.

\emph{Propositional half.}  Attach \texttt{thm:prop-lk-completeness} to the
theorem at 2488 and replace its statement (2488--2491) by: for all finite
multisets \(\Gamma,\Delta\) of propositional formulas, the sequent
\(\Gamma\Rightarrow\Delta\) is derivable in \(\mathsf{LK}_{\mathrm{CPL}}\)
iff every valuation making every formula of \(\operatorname{supp}(\Gamma)\)
true makes some formula of \(\operatorname{supp}(\Delta)\) true ---
word-for-word the sequent-satisfaction convention of 2156--2161.  Left to
right is \texttt{thm:prop-lk-soundness} (2295).  Right to left is the
reduction (1)--(5) verbatim with \texttt{thm:prop-completeness} (2433) in
place of \texttt{thm:fol-completeness}, the translation sketch at 2493--2507
as the retained subroutine, and the same initial sequents and structural
rules (2178--2234); no choice principle enters (the propositional Lindenbaum
lemma runs on a countable enumeration).  Landing: statement replaced at
2488--2491; steps (1)--(5) appended to the proof at 2493--2507.  Status
afterwards: likewise \textit{Status: proof sketch} while its translation
trees remain omitted.

\paragraph{Acceptance.}
(1) The completeness theorem at 4165 quantifies over sequents
\(\Gamma\Rightarrow\Delta\) with both contexts arbitrary finite multisets, and
its semantic condition is word-for-word the satisfaction convention of
3579--3581, matching the scope of \texttt{thm:fol-lk-soundness}.  (2) The proof
displays the reduction chain of steps (1)--(5), including the
\(\;\Rightarrow\delta_j,\neg\delta_j\) construction and the empty-succedent
case.  (3) The status label matches the evidence actually present: sketch
while the translation trees are omitted, complete only if they are supplied.
(4) The propositional theorem at 2488 likewise quantifies over sequents with
both contexts arbitrary finite multisets, matches the scope of
\texttt{thm:prop-lk-soundness}, is proved by the same reduction without any
choice principle, and carries the same status discipline; after both
restatements, soundness and completeness are paired over the same sequent
class in each calculus, discharging both line pairs of the register row.
\endgroup
